\scanpage{39}
\mspage{- 29 -}
% block: top; lines 1-3; display
\[
\begin{aligned}
\handwritten{\frac}{\handwritten{\partial}}{\handwritten{\partial}x_k}\left(\handwritten{\frac}{\handwritten{\partial}\unsure{u|n}_j}{\handwritten{\partial}x_i}\right)
&=\handwritten{\frac}{\handwritten{\partial}U_{ij}}{\handwritten{\partial}x_k}(\handwritten{x};\handwritten{u}(\handwritten{x}))+\handwritten{\sum_{h=1}^{q}}\handwritten{\frac}{\handwritten{\partial}U_{ij}}{\handwritten{\partial}y_h}(\handwritten{x};\handwritten{u}(\handwritten{x}))\,\handwritten{\cdot}\,\handwritten{\frac}{\handwritten{\partial}u_h}{\handwritten{\partial}x_k}\\
&=\handwritten{\frac}{\handwritten{\partial}U_{ij}}{\handwritten{\partial}x_k}(\handwritten{x};\handwritten{u}(\handwritten{x}))+\handwritten{\sum_{\unsure{k|h}=1}^{q}}\handwritten{\frac}{\handwritten{\partial}U_{ij}}{\handwritten{\partial}y_h}(\handwritten{x};\handwritten{u}(\handwritten{x}))U_{kh}(\handwritten{x};\handwritten{u}(\handwritten{x})).
\end{aligned}
\qquad\begin{gathered}
(1\handwritten{\leq}i\handwritten{\leq}p,\\
1\handwritten{\leq}k\handwritten{\leq}p,\\
1\handwritten{\leq}j\handwritten{\leq}q
\end{gathered}
\]
% The last bound meets the right image edge; no final parenthesis is visible.
% block: top-middle; lines 4-6
Ecriva-nt que les expressions ainsi obtenues pour $\handwritten{\frac}{\handwritten{\partial}}{\handwritten{\partial}x_k}(\handwritten{\frac}{\handwritten{\partial}u_j}{\handwritten{\partial}x_i})$ et $\handwritten{\frac}{\handwritten{\partial}}{\handwritten{\partial}x_i}(\handwritten{\frac}{\handwritten{\partial}u_j}{\handwritten{\partial}x_k})$ sont égales au point $\handwritten{x}$, puis faisant $\handwritten{x}=\handwritten{x}_0$ et tenant compte de (23) il vient
% block: middle; lines 7-8; display
\[
\begin{aligned}
&\handwritten{\frac}{\handwritten{\partial}U_{ij}}{\handwritten{\partial}x_k}(\handwritten{x}_0;\handwritten{y}_0)+\handwritten{\sum_{h=1}^{q}}\handwritten{\frac}{\handwritten{\partial}U_{ij}}{\handwritten{\partial}y_h}(\handwritten{x}_0;\handwritten{y}_0)U_{kh}(\handwritten{x}_0;\handwritten{y}_0)={}\\
&={}\handwritten{\frac}{\handwritten{\partial}U_{kj}}{\handwritten{\partial}x_i}(\handwritten{x}_0;\handwritten{y}_0)+\handwritten{\sum_{\unsure{k|h}=1}^{q}}\handwritten{\frac}{\handwritten{\partial}U_{kj}}{\handwritten{\partial}y_h}(\handwritten{x}_0;\handwritten{y}_0)U_{ih}(\handwritten{x}_0;\handwritten{y}_0)
\end{aligned}
\]
% block: middle; lines 9-10
Mais comme par hypothèse, $(\handwritten{x}_0,\handwritten{y}_0)$ est un point arbitraire de A, on obtient bien (24).

% block: middle; line 11
$2^{\handwritten{\mathrm{o}}}$ Les conditions sont \uline{suffisantes}. Considérons l'équation différentielle
% block: middle; line 12; display
\[
\handwritten{y}'=\underline{U}(\handwritten{x}_0+t(\handwritten{x}-\handwritten{x}_0);\handwritten{y})\handwritten{\cdot}(\handwritten{x}-\handwritten{x}_0)
\tag*{(25)}
\]
% block: middle-bottom; lines 13-19
où \uline{U} est la matrice $(U_{ij})$, $\handwritten{x}_0$ et $\handwritten{x}$ sont fixés dans B, et $\handwritten{y}'=\mathrm{d}\handwritten{y}/\mathrm{d}t$, t variant dans un voisina ge de 0 dans K. Il résulte de la formation de la série majorante dans la démonstration du th.2 que, dès que $\handwritten{x}-\handwritten{x}_0$ est assez petit, l'équation (25) admet une solution t une seule $\handwritten{v}(t;\handwritten{x}_0,\handwritten{x},\handwritten{y}_0)$ holomorphe pour $\handwritten{|}t\handwritten{|<}2$ et pour $\handwritten{x}$ et $\handwritten{x}_0$ dans un voisinage assez petit U de $\handwritten{a}$, $\handwritten{y}_0$ dans un voisinage assez petit V de $\handwritten{b}$ et telle que $\handwritten{v}(0;\handwritten{x}_0,\handwritten{x},\handwritten{y}_0)=\handwritten{y}_0$. Si $\handwritten{v}=(v_j)$, on a donc
% block: bottom; line 20; display
\[
\handwritten{\frac}{\mathrm{d}v_j}{\mathrm{d}t}=\handwritten{\sum_{i=1}^{p}}U_{ij}(\handwritten{x}_0+t(\handwritten{x}-\handwritten{x}_0);\handwritten{v}(t))(x_i-x_{i0})\qquad(1\handwritten{\leq}j\handwritten{\leq}q)
\tag*{(26)}
\]
% block: bottom; line 21
On déduit de là que
% block: bottom; lines 22-24; display
\[
\begin{aligned}
\handwritten{\frac}{\mathrm{d}}{\mathrm{d}t}\left(\handwritten{\frac}{\handwritten{\partial}v_j}{\handwritten{\partial}x_k}\right)&=\handwritten{\frac}{\handwritten{\partial}}{\handwritten{\partial}x_k}\left(\handwritten{\frac}{\mathrm{d}v_j}{\mathrm{d}t}\right)=U_{kj}(\handwritten{x}_0+t(\handwritten{x}-\handwritten{x}_0);\handwritten{v}(t))+\\
&\quad+t\handwritten{\sum_i}(x_i-x_{i0})\handwritten{\frac}{\handwritten{\partial}U_{ij}}{\handwritten{\partial}x_k}(\handwritten{x}_0+t(\handwritten{x}-\handwritten{x}_0);\handwritten{v}(t))\\
&\quad+\handwritten{\sum_{i,h}}(x_i-x_{i0})\handwritten{\frac}{\handwritten{\partial}U_{ij}}{\handwritten{\partial}y_h}(\handwritten{x}_0+t(\handwritten{x}-\handwritten{x}_0);\handwritten{v}(t))\handwritten{\frac}{\handwritten{\partial}v_h}{\handwritten{\partial}x_k}
\end{aligned}
\qquad(1\handwritten{\leq}j\handwritten{\leq}q,\ 1\handwritten{\leq}k\handwritten{\leq}p)
\]
% archive sheet number: 38
